% Standalone problem document, generated from the adjacent README.md.
% Compile with XeLaTeX. Mathematical content is maintained in Markdown.
\documentclass[11pt,a4paper]{article}
\usepackage[fontsize=10.5pt]{scrextend}
\usepackage[margin=22mm,headheight=15pt,headsep=9mm,footskip=11mm]{geometry}
\usepackage{fontspec}
\setmainfont{texgyrepagella}[Extension=.otf,UprightFont=*-regular,BoldFont=*-bold,ItalicFont=*-italic,BoldItalicFont=*-bolditalic]
\setsansfont{texgyreheros}[Extension=.otf,UprightFont=*-regular,BoldFont=*-bold,ItalicFont=*-italic,BoldItalicFont=*-bolditalic]
\setmonofont{texgyrecursor}[Extension=.otf,UprightFont=*-regular,BoldFont=*-bold,ItalicFont=*-italic,BoldItalicFont=*-bolditalic,Scale=MatchLowercase]
\usepackage{amsmath,amssymb}
\usepackage{unicode-math}
\setmathfont{texgyrepagella-math.otf}
\usepackage{xcolor,microtype,enumitem,needspace,fancyhdr}
\usepackage{bookmark}
\definecolor{ink}{HTML}{17344B}
\definecolor{accent}{HTML}{197D80}
\definecolor{muted}{HTML}{596773}
\hypersetup{colorlinks=true,linkcolor=accent,urlcolor=accent,pdftitle={KE-04 - Strict
interlacing across block Lanczos
iterations},pdfauthor={Open Problems in Numerical Linear Algebra}}
\urlstyle{same}
\setlength{\parindent}{0pt}
\setlength{\parskip}{5pt plus 1pt minus 1pt}
\linespread{1.04}
\setlength{\emergencystretch}{3em}
\widowpenalty=10000
\clubpenalty=10000
\displaywidowpenalty=10000
\allowdisplaybreaks[1]
\setlist{nosep,leftmargin=1.5em}
\providecommand{\tightlist}{\setlength{\itemsep}{0pt}\setlength{\parskip}{0pt}}
\providecommand{\pandocbounded}[1]{#1}
\pagestyle{fancy}
\fancyhf{}
\fancyhead[L]{\sffamily\footnotesize\color{muted}OPEN PROBLEMS IN NUMERICAL LINEAR ALGEBRA}
\fancyhead[R]{\sffamily\bfseries\footnotesize\color{accent}KE-04}
\fancyfoot[L]{\parbox[b]{0.9\textwidth}{\sffamily\fontsize{8}{10}\selectfont\color{muted}Editorial ratings; a bounded search cannot certify that a problem remains unsolved.\\Literature check: 2026-09-11}}
\fancyfoot[R]{\sffamily\footnotesize\color{muted}\thepage}
\renewcommand{\headrulewidth}{0.3pt}
\renewcommand{\headrule}{\hbox to\headwidth{\color{accent}\leaders\hrule height \headrulewidth\hfill}}
\makeatletter
\renewcommand{\section}{\Needspace{5\baselineskip}\@startsection{section}{1}{0pt}{12pt plus 2pt minus 2pt}{4pt}{\sffamily\bfseries\large\color{ink}}}
\renewcommand{\subsection}{\Needspace{4\baselineskip}\@startsection{subsection}{2}{0pt}{9pt plus 1pt minus 1pt}{3pt}{\sffamily\bfseries\normalsize\color{ink}}}
\makeatother
\setcounter{secnumdepth}{0}
\begin{document}
{\sffamily\small\color{muted}Eigenvalues and inverse problems\par}
\vspace{3pt}
{\raggedright\hyphenpenalty=10000\exhyphenpenalty=10000\sffamily\bfseries\fontsize{21}{24}\selectfont\color{ink}Strict
interlacing across block Lanczos iterations\par}
\vspace{7pt}
{\sffamily\small\textbf{Difficulty:} challenging\quad\textbf{Importance:} interesting
to the community\par}
{\sffamily\small\bfseries\color{accent}Status: Solution claimed\par}
\vspace{4pt}
{\color{accent}\hrule height 0.8pt}
\vspace{6pt}
\textbf{Rating rationale:} Challenging reflects a stronger strict
interlacing law than general compression interlacing supplies; community
impact is spectral information across block Krylov iterations.

\subsection{Resolution claim ---
2026-09-11}\label{resolution-claim-2026-09-11}

\textbf{Affirmative resolution claimed; submitted by Matthew Colbrook.}
See the
\href{https://github.com/ajt60gaibb/OpenProblemsInNLA/blob/main/eigenvalues-and-inverse-problems/KE-04/solution.md}{complete
manuscript}, \textbf{Theorem KE-04, sections 1--3}
(\href{https://github.com/ajt60gaibb/OpenProblemsInNLA/blob/main/eigenvalues-and-inverse-problems/KE-04/solution.pdf}{PDF}
·
\href{https://github.com/ajt60gaibb/OpenProblemsInNLA/blob/main/eigenvalues-and-inverse-problems/KE-04/solution.tex}{LaTeX}),
prepared 11 September 2026.

The manuscript claims strict interval occupancy for every allowed pair
of block Lanczos iterations and every indicated index, in exact
arithmetic before the first loss of full block dimension. Its
quadratic-polynomial argument includes eigenvalue multiplicities and
excludes coincident interval endpoints in the stated range.

The manuscript discloses generation in a ChatGPT conversation. Its full
proof has not been independently verified here, and no publication or
priority claim is made.
\href{https://github.com/ajt60gaibb/OpenProblemsInNLA/blob/main/references/colbrook-2026-09-11/README.md}{Supporting
diagnostics} were rerun successfully; these are not an independent proof
audit. This update checks the submitted claim against the displayed
target; the dated literature checks below remain the record of the
earlier search. The ratings above are historical, and this entry is
excluded from the open count.

\subsection{Original problem
statement}\label{original-problem-statement}

Let \(A\in\mathbb R^{n\times n}\) be symmetric and
\(V\in\mathbb R^{n\times p}\) have full column rank. Write \[
\mathcal K_j=\operatorname{range}[V,AV,\ldots,A^{j-1}V],
\] and let \(s\) be the largest integer with \(\dim\mathcal K_s=sp\).
For \(1\le j\le s\), choose an orthonormal basis \(Q_j\) of
\(\mathcal K_j\) and order the eigenvalues of \(T_j=Q_j^TAQ_j\) as
\(\theta_1^{(j)}\le\cdots\le\theta_{jp}^{(j)}\), with multiplicity.

For every \(1\le k<j\le s\) and \(1\le i\le(k-1)p\), must the open
interval \[
(\theta_i^{(k)},\theta_{i+p}^{(k)})
\] contain an eigenvalue of \(T_j\)? The statement concerns exact
arithmetic before the first loss of full block dimension. Although
Lanczos supplies compatible block tridiagonal representations, the
spectra do not depend on the chosen bases.

This would extend a scalar Lanczos property used to interpret later Ritz
values and distinguish genuine eigenvalue approximation from clusters
created by finite precision.

\newpage
\subsection{References}\label{references}

D. Šimonová and P. Tichý, \emph{On finite precision block Lanczos
computations}, arXiv:2507.16484v1 (22 July 2025), §2.1, unnumbered
conjecture and final discussion
(\href{https://arxiv.org/html/2507.16484v1}{primary manuscript}). J.
Liesen and Z. Strakoš, \emph{Krylov Subspace Methods: Principles and
Analysis}, Oxford 2013, the scalar Lanczos/orthogonal-polynomial
background cited by the source
(\href{https://doi.org/10.1093/acprof:oso/9780199655410.001.0001}{publisher}).

\subsection{Status check --- 2026-09-08}\label{status-check-2026-09-08}

The arXiv abstract/history and full §2.1 were inspected; v1 remains the
only listed version. Searches ``block Lanczos interlacing conjecture'',
``Šimonová Tichý interlacing'', and 2026 proof/counterexample queries
found no resolution. Standard weak Cauchy interlacing is not the
asserted strict interval-occupancy result.

\subsection{Audit update --- 2026-09-10}\label{audit-update-2026-09-10}

Rechecked the explicit conjecture and concluding discussion in the
\href{https://arxiv.org/html/2507.16484v1}{2025 block-Lanczos
manuscript}. Its block-width index ranges agree with this entry.
Searches for later strict block-Lanczos interlacing proofs found no
resolution; ordinary Cauchy interlacing is weaker than the displayed
claim.
\end{document}
